\section{G-injectivity under physical maps}%
\label{sec:peps_g_injective_physical_map}

\begin{theorem}[Physical maps preserving G-injectivity]%
    \label{thm:peps_g_injective_physical_map}
    \lean{TNLean.PEPS.comp_invariant_of_invariant,
        TNLean.PEPS.IsGInjective.comp_of_injOn_range,
        TNLean.PEPS.IsGInjective.comp_of_injective,
        TNLean.PEPS.IsGInjective.comp_equiv,
        TNLean.PEPS.isGInjective_comp_equiv_iff,
        TNLean.PEPS.IsGInjective.comp_iff_injOn_range,
        TNLean.PEPS.IsGInjective.comp_iff_range_inf_ker_eq_bot,
        TNLean.PEPS.IsGInjective.comp_iff_injective_rangeSubtype}
    \leanok
    \uses{def:peps_g_injective}
    Let $T:\mathcal V\to\mathcal P$ be $G$-injective and let
    $F:\mathcal P\to\mathcal Q$ be linear.
    The composite $FT$ is invariant under the original virtual
    representation. If $F|_{\operatorname{im}T}$ is injective, then
    $FT$ is $G$-injective. For a finite group $G$, the converse holds:
    $FT$ is $G$-injective if and only if
    $\operatorname{im}T\cap\ker F=\{0\}$.
    In particular an invertible physical change of coordinates
    preserves $G$-injectivity in both directions.
    These are local consequences of
    \cite[Definition~5.1]{Schuch2010PEPS}.
\end{theorem}

\begin{proof}\leanok
    \uses{thm:peps_g_injective_invariants_range}
    Invariance follows from $FTU_g=FT$.
    If $x\in\mathcal V^G$ and $FTx=0$, injectivity of $F$ on
    $\operatorname{im}T$ gives $Tx=0$, and $G$-injectivity gives
    $x=0$. Conversely, for finite $G$, every
    $y\in\operatorname{im}T$ is $Tx$ for an invariant vector $x$.
    If $Fy=0$ and $FT$ is $G$-injective, then $x=0$, hence $y=0$.
    For an invertible physical map, apply the same implication to its
    inverse to obtain the reverse direction.
\end{proof}

\begin{theorem}[Physical isometries preserving G-isometry]%
    \label{thm:peps_g_isometric_physical_map}
    \lean{TNLean.PEPS.IsGIsometric.comp_of_inner_eq,
        TNLean.PEPS.IsGIsometric.comp_isometry,
        TNLean.PEPS.isGIsometric_comp_equiv_iff}
    \leanok
    \uses{def:peps_g_injective}
    Let $T$ be $G$-isometric with normalization factor $c>0$.
    A physical linear isometry $F$ makes $FT$ $G$-isometric with the
    same factor $c$. This includes rectangular isometries.
    An invertible physical isometry preserves $G$-isometry in both
    directions.
\end{theorem}

\begin{proof}\leanok
    \uses{thm:peps_g_injective_physical_map}
    An isometry is injective, so $FT$ is $G$-injective.
    For invariant $x,y$,
    $\langle FTx,FTy\rangle=\langle Tx,Ty\rangle
        =c\langle x,y\rangle$.
    For a matrix isometry, the first equality follows from
    $F^\dagger F=\Id$. The inverse of an invertible isometry is an
    isometry, giving the converse.
\end{proof}

\begin{theorem}[Compatibility with virtual comparison]%
    \label{thm:peps_g_injective_physical_virtual_comparison}
    \lean{TNLean.PEPS.IsGInjective.virtualEquiv_comp_of_injective}
    \leanok
    \uses{def:peps_g_injective_virtual_comparison}
    Let $T:\mathcal V\to\mathcal P$ and
    $S:\mathcal W\to\mathcal P$ be injective on the invariant
    spaces of possibly different finite groups, and suppose their
    physical ranges agree. Applying the same injective physical map
    $F:\mathcal P\to\mathcal Q$ to both tensors leaves their
    virtual comparison unchanged.
\end{theorem}

\begin{proof}\leanok
    \uses{thm:peps_g_injective_virtual_comparison}
    The comparison for $FT$ and $FS$ satisfies $FS\Phi x=FTx$.
    Injectivity of $F$ gives $S\Phi x=Tx$, which uniquely
    characterizes the original comparison.
\end{proof}

\begin{definition}[Physical deformation of a regular site]%
    \label{def:peps_regular_physical_site}
    \lean{TNLean.PEPS.regularPhysicalMapSite,
        TNLean.PEPS.regularSiteMap_regularPhysicalMapSite}
    \leanok
    \uses{def:peps_regular_site_gram}
    For a site tensor $a(\eta,s)$ whose virtual labels are group
    elements on an arbitrary incident-leg set, define
    $a^F(\eta,t)=\sum_sF_{ts}a(\eta,s)$.
    Its physical coefficient map is $T_{a^F}=FT_a$.
\end{definition}

\begin{theorem}[Physical deformation of G-injective site families]%
    \label{thm:peps_g_injective_physical_site_family}
    \lean{TNLean.PEPS.IsGInjective.physicalMapSite,
        TNLean.PEPS.IsGIsometric.physicalMapSite,
        TNLean.PEPS.IsGInjective.regularPhysicalMapSite,
        TNLean.PEPS.IsGIsometric.regularPhysicalMapSite,
        TNLean.PEPS.physicalDeform_groupBondTensor,
        TNLean.PEPS.isGInjective_regularPhysicalMapSite_family,
        TNLean.PEPS.isGIsometric_regularPhysicalMapSite_family}
    \leanok
    \uses{def:peps_regular_physical_site,def:peps_physical_deformation}
    An injective physical map preserves $G$-injectivity of a site
    tensor, and an isometry preserves its $G$-isometry.
    Both statements hold for four virtual legs and for arbitrary
    incident-leg sets.
    Consequently invertible physical maps at every vertex preserve
    regular $G$-injectivity of a complete finite-graph site family;
    physical isometries preserve regular $G$-isometry.
    In finite bond coordinates the resulting tensor family is exactly
    the physical deformation of the original group-labelled family.
\end{theorem}

\begin{proof}\leanok
    \uses{thm:peps_g_injective_physical_map,thm:peps_g_isometric_physical_map}
    Apply the preceding composition results to each site's coefficient
    map. The group-labelled coordinates identify the deformed
    coefficients with the same physical sums in finite bond coordinates.
\end{proof}
